\chapter{绪论}
\label{ch:introduction}

\section{研究背景及意义}
\label{sec:1_1}

随着低轨卫星通信系统的快速发展，空天地一体化网络正在成为新一代通信体系的重要组成部分。以 Starlink 为代表的低轨卫星星座在推动广域接入和应急通信能力提升的同时，也带来了新的监管问题\cite{Stock2025LEOPNTSurvey,Fan2024MassiveSoOP,Kozhaya2025StarlinkPNT}。非合作低轨卫星终端可能通过异常接入、违规辐射或非法占用链路资源，对无线电管理秩序和通信安全造成潜在影响。因此，面向非合作低轨卫星终端的快速发现与定位，已经成为复杂电磁环境下值得关注的现实问题\cite{Yin2025LEO}。

与常规通信终端不同，非合作低轨卫星终端通常不主动提供同步信息和位置辅助信息，监测系统只能依赖其辐射信号及相关传播特征进行终端感知与定位\cite{Stock2025LEOPNTSurvey,Humphreys2023StarlinkSignal,Torrieri1984PassiveLocation,Weiss2011DirectGeolocation}。与此同时，这类目标往往分布在高层建筑密集区域、街道走廊或复杂城市边缘地带，传播过程中普遍存在遮挡、反射、绕射和散射等现象。受复杂城市环境影响，传统定位方法所依赖的理想视距条件往往难以满足，单一观测量容易出现失真，定位精度和稳定性均受到明显制约\cite{TR38901,Guvenc2009TOASurvey,Wang2013TDOANLOS,deSousa2018NLOSRayTracing}。

在此背景下，无人机平台为非合作终端排查提供了新的实现路径。相较于固定监测站，无人机具有机动性强、部署灵活和空间视角可调等特点，能够在较短时间内进入目标区域上空或周边，对地面终端开展机动观测。尤其在复杂街区和局部遮挡场景下，无人机更容易形成有效观测几何，因此在非合作终端排查任务中具有较高的应用潜力\cite{AlHourani2014A2GPathLoss,Chen2017UAVRadioMap,Li2022UAVRadioMapSemantics,SpectrumSurveying2022}。与此同时，低成本无人机平台在载荷、续航、算力和阵列规模等方面又存在明显约束，这决定了相关定位方法不能简单照搬高成本、大规模地面监测系统的技术路线，而需要围绕平台特点重新设计\cite{Wymeersch2009CooperativeLocalization}。

从现有方法看，复杂城市环境下机载平台的非合作终端定位仍面临两方面突出问题。其一，单目标场景下，不同观测量对环境变化和运动状态的敏感性并不一致，其有效性会随传播条件和轨迹过程持续变化，导致固定权重或静态融合规则难以稳定适应复杂场景。其二，多目标场景下，无人机空间扫描所形成的观测通常具有稀疏性和不均匀性，这使得目标区域的空间响应表达不完整，容易引发局部多峰和多径鬼影问题；与此同时，多无人机协同过程中还可能出现节点掉点现象，从而进一步加大定位难度\cite{RadioUNet2021,Romero2022SpectrumCartography}。如何结合无人机平台的机动观测特点，在复杂环境中实现稳定、准确且具有鲁棒性的定位识别，是本文研究的核心问题。

本文围绕低成本无人机平台下非合作低轨卫星终端的定位识别问题开展研究，分别面向单目标连续定位与多目标并发定位两类任务展开方法设计。前者重点解决复杂环境中多源观测的自适应融合问题，后者重点解决空间扫描条件下的目标识别、信号成像和鲁棒协同定位问题。相关研究不仅面向非合作终端排查这一具体应用，也对复杂环境中的无人机无线感知与智能定位具有一定参考价值。

\section{国内外研究现状}
\label{sec:1_2}

\subsection{定位技术研究现状}
\label{subsec:1_2_1}

随着低轨卫星通信系统的快速发展，围绕低轨机会信号利用、复杂环境传播建模以及无人机机动感知的研究不断增多。整体来看，相关工作正在由传统地面固定站条件下的几何定位，逐步转向低轨卫星信号利用、复杂城市传播机理分析以及机动平台协同感知等方向\cite{Stock2025LEOPNTSurvey,Fan2024MassiveSoOP,Yin2025LEO}。

在低轨卫星信号底层结构解析方面，相关研究已经从早期的可行性验证逐步发展到面向定位的结构化处理。Humphreys等人分别从 Starlink Ku 波段下行链路的信号结构解析和载波相位跟踪两个角度开展研究，揭示了其同步结构、参考特征以及可用于定位处理的底层物理线索\cite{Humphreys2023StarlinkSignal,Humphreys2023StarlinkCarrier}。随后，Dureppagari等人进一步指出，Starlink 下行链路中存在类正交频分复用（Orthogonal Frequency Division Multiplexing, OFDM）的可利用结构，为更稳定的相关积累和精确定位处理提供了依据\cite{Dureppagari2024StarlinkOFDM}。Kozhaya 等人的工作则进一步从定位、导航与授时（Positioning, Navigation and Timing, PNT）角度系统总结了 Starlink 机会信号在观测量提取与定位实验中的可行链路\cite{Kozhaya2025StarlinkPNT}。这些研究表明，非合作条件下的低轨卫星下行信号并非完全不可利用，其底层结构具有被逆向解析并服务于定位感知的现实可能性。

然而，低轨卫星下行信号的可利用性并不意味着可直接获得高精度测距结果。Qin等人进一步分析了 Starlink Ku 波段下行链路的定时特性，指出其在短时尺度上存在明显的非连续跃变和相位抖动\cite{Qin2025Timing}。这说明低轨卫星下行信号虽然可为定位提供一定观测基础，但其时间稳定性并不完全满足传统高精度到达时间（Time of Arrival, TOA）或到达时间差（Time Difference of Arrival, TDOA）定位对严格同步的要求，因此直接将其视为理想测距基准仍存在明显局限。从无线定位与终端排查的系统应用角度看，Yin等人\cite{Yin2025LEO}系统梳理了国内相关研究进展，指出低轨卫星终端监测技术已经由早期的单站监测逐步走向多平台协同监测。

在复杂城市环境中，传统几何定位技术仍是研究的重要基础，但其局限也越来越清晰。3GPP 技术报告（Technical Report, TR）38.901 给出了统一的城市传播场景建模框架，强调多径簇、时延扩展、角扩展以及视距/非视距（Line-of-Sight / Non-Line-of-Sight, LOS/NLOS）条件对无线观测结果的系统性影响\cite{TR38901}。针对空地链路和卫星地面链路，Al-Hourani 等人以及 Gapeyenko 等人分别从低空平台路径损耗和卫星到地面链路衰落的角度说明，仰角、遮挡和城市几何结构会显著改变观测质量分布\cite{AlHourani2014A2GPathLoss,Gapeyenko2017SatGroundPathLoss}。这些研究从传播机理层面表明，复杂城市环境中的定位误差并不只是随机噪声放大，而是与遮挡条件、多径结构和观测几何共同耦合的系统性退化。

在算法层面，Chan 和 Ho 的双曲定位估计器给出了经典 TDOA 定位的高效求解框架，而后续关于 TOA/TDOA 与 NLOS 偏差的综述工作则系统总结了复杂环境下的典型退化机理与缓解路径\cite{Chan1994Hyperbolic,Guvenc2009TOASurvey}。进一步地，Wang和Ho针对 NLOS 条件下的 TDOA 偏差提出了鲁棒估计方法，Dickerson等人则基于真实无人机飞行数据分析了 TDOA 定位在城市混合视距与非视距环境中的性能变化，结果表明受非视距偏差、观测几何和带宽条件共同影响，TDOA 误差在复杂场景中会出现显著发散\cite{Wang2013TDOANLOS,Dickerson2025TDOA,deSousa2018NLOSRayTracing}。这些研究共同说明，在高层建筑密集区域，仅依赖单一几何观测建立位置约束往往难以获得稳定结果，复杂环境下的终端定位必须引入更强的环境感知和误差补偿机制。

除几何观测外，围绕无线电环境重构与空间信号场建模的研究近年来发展较快。Zappone等人和 Jiang等人分别从环境感知通信和信道知识地图（Channel Knowledge Map, CKM）的角度，提出将无线传播环境组织为可查询、可更新的无线电地图（Radio Map）/无线电环境地图（Radio Environment Map, REM）中间表示，用于刻画目标区域内信号强度的空间分布结构，并为后续定位和感知任务提供统一表征\cite{Zappone2021CKM,Jiang2022CKMTutorial}。在此基础上，Li等人将城市环境语义显式引入 UAV 辅助 Radio Map 构建过程，提高了重构精度与可解释性\cite{Li2022UAVRadioMapSemantics}；Bazerque和Giannakis则通过贝叶斯主动学习说明了采样策略与 Radio Map 重构质量之间的强耦合关系\cite{Bazerque2010BayesianRadioMap}。这些工作表明，Radio Map/REM 已从早期的经验插值工具逐步演化为融合环境语义、传播机理与空间结构先验的中间表示层。

随着深度学习方法引入频谱制图领域，RadioUNet、深度补全式 Radio Map 估计以及数据驱动的频谱制图方法进一步表明，空间信号场重构已经从传统插值逐步发展为融合空间结构先验和学习表征能力的系统问题\cite{RadioUNet2021,Teganya2022DeepCompletion,Romero2022SpectrumCartography,Feng2025RadioMapSurvey}。进一步地，Ye等人提出物理知识启发扩散模型以改善稀疏采样条件下的 REM 重构质量，HiFiRadio 相关研究和 RadioTransformer 模型则分别从高保真三维重建和多尺度特征表达角度推动了复杂环境中无线电地图构建技术的发展\cite{Ye2025PhysicsREM,HiFiRadio2025,RadioTransformer2025}。与此相对应，多阶段无人机采样与航迹规划方法以及自主 UAV 频谱测绘工作表明，对于低成本无人机平台而言，定位性能不仅取决于后端算法，也与前端空间观测策略密切相关\cite{UAVREMPlanning2025,SpectrumSurveying2022,Chen2017UAVRadioMap}。

总体来看，国内外定位技术研究已经从传统接收信号强度指示（Received Signal Strength Indicator, RSSI）、到达角（Angle of Arrival, AoA）、到达时间（Time of Arrival, TOA）和到达时间差（Time Difference of Arrival, TDOA）等经典观测量，逐步扩展到低轨机会信号利用、复杂城市传播建模、无线电环境地图重构以及无人机机动协同感知等方向。现有研究为复杂场景中的非法低轨卫星终端排查提供了重要基础，但仍存在一些不足：其一，低轨卫星下行信号的底层可利用特征虽已逐步明确，但与非法终端排查任务直接衔接的技术链路仍不完整；其二，传统几何定位在复杂城市环境中的退化机理虽然已较清楚，但针对低成本无人机平台约束条件的适配方案仍显不足；其三，空间信号场重构与目标定位识别之间尚缺乏更紧密的任务耦合。因此，围绕低成本无人机平台开展复杂环境下非法低轨卫星终端定位研究，仍具有较强的现实意义。

\subsection{定位算法研究现状}
\label{subsec:1_2_2}

如果说前述研究更多关注低轨信号能否被利用以及复杂环境对观测模型会造成怎样的影响，那么定位算法研究则主要聚焦于如何在这些约束条件下实现稳定、准确的位置推断。总体来看，当前相关算法研究大致可以分为三类：一类是基于传统几何测量与滤波估计的方法，一类是基于深度学习的表征与检测方法，另一类是面向动态环境与多节点协同的强化学习和知识迁移方法。

在传统方法方面，基于 TDOA、AoA 和滤波估计的联合定位仍然是当前研究中的重要基线。相关研究中，已有学者提出基于扩展卡尔曼滤波的 TDOA 与 AoA 融合定位方法\cite{EKFTDOAAOA2025}，将方向观测和时间差观测统一纳入状态估计框架，在一定程度上缓解了单一观测量的不稳定问题；同时，经典双曲定位估计器和面向 NLOS 偏差的鲁棒变换与优化方法也被用于提升复杂环境中的定位稳定性\cite{Chan1994Hyperbolic,Wang2013TDOANLOS}。这类方法的优势在于物理意义清晰、求解路径成熟，适合作为复杂环境定位的参考基线。但从已有研究结果来看，这类方法仍然较强依赖于噪声统计和观测模型的正确性，当非视距偏差、多径扰动和异常观测显著增多时，其性能仍会明显下降。也就是说，传统联合滤波能够改善单源观测条件下的定位效果，但尚不足以从根本上解决复杂环境中观测质量持续波动的问题。

随着深度学习的发展，研究者开始将空间信号场重构、环境表征学习与目标区域判别统一到数据驱动框架中。相关工作表明，基于高保真 Radio Map/REM 的中间表示能够在复杂场景下保留更丰富的空间纹理和传播环境信息，从而为后续定位识别提供更具判别性的输入\cite{HiFiRadio2025,RadioTransformer2025}。进一步地，环境语义感知的 UAV 辅助 Radio Map 构建方法\cite{Li2022UAVRadioMapSemantics}以及物理知识驱动的 REM 重构方法\cite{Ye2025PhysicsREM}说明，仅依赖传统数值插值已难以满足复杂城市环境中定位识别任务对精度与泛化能力的双重要求。相比于传统几何定位，这类方法更强调利用空间场分布、区域纹理和上下文结构完成目标识别与定位估计。

除静态学习方法外，面向动态环境与观测质量持续变化的强化学习方法也受到广泛关注。Mnih 等人提出的深度 Q 网络（Deep Q-Network, DQN）\cite{Mnih2015DQN}以及 Lillicrap 等人提出的深度确定性策略梯度方法\cite{ddpg}为复杂状态下的自适应决策提供了通用框架；在定位相关应用中，深度强化学习已被用于协同定位调度、无人机干扰源搜索和多无人机协同决策等任务\cite{Peng2019DRLLocalization,Wu2019QLearningLocalization,MultiUAVDRL2025,Drones2025CooperativeTarget}。这类方法的共同特点在于，将观测选择、权重调整或协同策略显式建模为时序决策过程，从而提高系统在动态环境中的自适应能力。

与此同时，知识迁移和知识蒸馏也为复杂模型向轻量模型迁移判别能力提供了可行路径。Hinton 等人提出的知识蒸馏框架\cite{Hinton2015KD}表明，教师模型中蕴含的判别知识可以通过软标签或中间表征迁移到学生模型中，从而缓解边缘设备算力受限和观测缺失带来的性能退化问题；而面向多智能体学习的蒸馏方法也说明，知识蒸馏不仅可以用于压缩，还可以服务于协同策略迁移与鲁棒决策\cite{Tseng2022OfflineMARLKD}。然而，现有研究大多面向一般无线感知、协同控制或模型压缩任务，针对低成本无人机平台下非法低轨卫星终端排查这一特定场景，仍缺少与传播机理、空间采样和多目标识别紧耦合的系统性算法设计。

\section{现有研究存在的主要问题}
\label{sec:1_3}
综上所述，国内外学者在低轨卫星信号解析、机会信号定位以及多传感器融合等方面已取得了一系列阶段性成果。然而，面向低成本无人机平台在复杂城市环境下排查非法卫星终端的实际任务，仍存在以下主要问题：
\begin{enumerate}
    \item 非合作信号时钟抖动制约定位鲁棒性：现有高精度定位方法（如精密 TDOA）高度依赖毫秒级以下的时间同步。然而，低轨卫星非法终端具有非合作特征，且其下行信号存在由软件调控引起的不规则跳变和时钟漂移 \cite{Qin2025Timing}，导致传统几何定位模型在面临此类非理想特性时容易产生较大的定位残差。
    \item 复杂传播环境导致单一观测量可靠性失稳：在城市峡谷等典型场景下，多径干扰和非视距（NLOS）传播会造成接收信号强度指示（Received Signal Strength Indicator, RSSI）、到达角（Angle of Arrival, AoA）和到达时间差（Time Difference of Arrival, TDOA）等观测量的测量噪声呈现显著的非高斯特征 \cite{Wang2013TDOANLOS,Dickerson2025TDOA}。现有研究多侧重于理想噪声模型，在动态多变的环境中，固定权的静态融合方法难以适应观测质量的剧烈波动。
    \item 稀疏采样条件下多目标识别面临鬼影干扰：利用无人机进行空间扫描时，获取的信号强度或角度特征通常具有空间稀疏性。在多目标并发环境下，多径反射产生的伪峰与真实目标容易混淆，形成“多径鬼影”现象。现有基于简单峰值检测的方法在低信噪比和强多径背景下识别准确率较低，缺乏有效的跨模态特征辅助识别机制。
\end{enumerate}

\section{本文主要研究内容}
\label{sec:1_4}
本文以低成本无人机为观测平台，针对复杂城市环境下非法低轨卫星终端的定位识别问题，重点围绕传播机理建模、多源智能融合以及多目标信号成像等方面展开深入研究。主要内容安排如下：
\begin{enumerate}
    \item 低成本无人机定位系统建模与传播机理分析：首先构建无人机机动观测下的非合作定位系统总体架构。针对复杂城市场景，建立基于射线追踪技术的低轨卫星信号传播模型，深入分析多径对 RSSI、AoA、TOA 等关键观测量的影响机理，为后续算法设计提供物理支撑。
    \item 基于物理感知深度强化学习的多源自适应融合定位方法：针对单目标定位中观测质量随环境波动的特点，研究基于深度强化学习的自适应融合框架。提出一种融合物理先验知识的物理感知深度 Q 网络（Physics-Aware Deep Q-Network, PA-DQN）智能体，通过建立物理感知的状态空间描述和奖励函数设计，实现对多源异构观测权重的实时自适应调节，提升复杂环境下的定位鲁棒性。
    \item 基于信号成像与跨模态检测的多目标并发定位方法：针对多目标场景下的识别与定位瓶颈，研究基于无线电地图（Radio Map）的信号成像重构方法。结合亥姆霍兹方程引入物理约束，实现稀疏采样下的高保真 REM 重构；并利用知识蒸馏提高掉点情况下的定位性能。提高并发多目标的识别准确率。
\end{enumerate}


\section{论文结构安排}
\label{sec:1_5}
本文共分为五章，各章节的具体内容安排如下：
\begin{itemize}
    \item 第 1 章：绪论。介绍研究背景及意义，梳理国内外研究现状，明确现有研究问题，并阐述本文的主要研究内容与创新点。
    \item 第 2 章：系统模型与定位理论基础。建立低成本无人机定位系统架构，分析复杂城市环境下的卫星信号传播机理，并介绍经典定位方法与本文所需的基础理论。
    \item 第 3 章：基于物理感知深度强化学习的单目标自适应融合定位方法。针对多源观测质量波动问题，提出 PA-DQN 融合框架，并进行实验验证与性能分析。
    \item 第 4 章：基于信号成像与知识蒸馏方法，研究掉点情况下基于 Radio Map 的定位算法与视觉驱动的多目标识别定位算法。
    \item 第 5 章：结论与展望。总结全文研究工作，指出创新点与不足，并对未来研究方向进行展望。
\end{itemize}
